\chapter{Launching}

La fase di Launching segna l'avvio ufficiale dell'esecuzione. In questa fase il team è stato messo nelle condizioni di lavorare in modo efficace, chiarendo ruoli e responsabilità, definendo le regole operative e stabilendo i canali di comunicazione e le procedure per la gestione dei cambiamenti.

\section{Project Kick-Off Meeting}

Il progetto è stato avviato ufficialmente con un Project Kick-Off Meeting tenutosi il 15 ottobre 2025, a cui hanno partecipato il team di PlayHeritage Labs, il committente Giovanni Marchetti in qualità di sponsor e l'esperta di dominio Francesca Giuliani. Durante l'incontro sono stati condivisi la visione e gli obiettivi del progetto, il perimetro dell'MVP, la timeline con le milestone critiche, il budget, i rischi principali e le regole operative, in modo da allineare tutti i partecipanti prima dell'inizio dello sviluppo. Il verbale completo, con l'agenda, le decisioni approvate e gli action items, è riportato nell'\allegatolink{Documentazione/Launching/Allegato4.1-ProjectKickOffMeeting.pdf}{Allegato 4.1 - Project Kick-Off Meeting}.

\section{RASCI Matrix}

Per assicurare il successo del progetto è stata redatta una matrice RASCI, che definisce per ogni attività chi è Responsible (esegue), Accountable (approva ed è unico responsabile del risultato), Support (fornisce supporto), Consulted (viene consultato per competenza) e Informed (viene tenuto aggiornato). La matrice assegna le responsabilità a livello di attività della WBS: i task di dettaglio ereditano la riga dell'attività di appartenenza, mentre le attività con una distribuzione di ruoli diversa --- come la validazione delle regole con l'esperta di dominio o la User Acceptance Testing --- hanno una riga dedicata. Copre tutti i sottosistemi e le attività di Project Management, garantendo che ogni attività abbia esattamente un Accountable ed eliminando così le ambiguità sui ruoli. La RASCI è riportata nell'\textbf{Allegato 4.2 - RASCI Matrix}, disponibile nella \allegatolink{Documentazione/Launching/Allegato4.2-RASCI-Visuale.pdf}{versione visiva a colori} e nella \allegatolink{Documentazione/Launching/Allegato4.2-RASCI.pdf}{versione testuale completa}.

\section{Regole Operative}

Per una buona riuscita del progetto sono state definite delle regole operative condivise da tutte le persone coinvolte. Esse comprendono una strategia di problem solving strutturata (basata sull'analisi della causa root con la tecnica dei ``5 Whys''), un framework di decision making articolato su tre livelli (operativo, tattico e strategico) secondo il principio ``decidere al livello più basso possibile'', un processo di conflict resolution progressivo in tre fasi (discussione diretta, mediazione facilitata e decisione esecutiva), le tecniche di brainstorming adottate e la definizione delle cerimonie di team (Daily Standup, Sprint Planning, Backlog Refinement, Sprint Review, Retrospective e Project Status Meeting). Le regole operative sono documentate nell'\allegatolink{Documentazione/Launching/Allegato4.3-RegoleOperative.pdf}{Allegato 4.3 - Regole Operative}.

\section{Comunicazione}

La comunicazione all'interno del team è considerata essenziale per la coesione e per la riduzione dei malintesi, ed è regolata da tre principi: comunicazione asincrona prima di tutto, trasparenza delle decisioni e tempi di risposta proporzionati alla priorità. Il canale primario è Slack, organizzato in canali dedicati --- comunicazioni generali, daily standup, discussioni tecniche e urgenze, queste ultime da leggere entro poche ore; l'email è riservata alle comunicazioni formali verso lo sponsor e gli enti esterni (report mensili, approvazioni di milestone, Change Request, contratti), mentre Notion funge da single source of truth: raccoglie le note dei meeting, gli status report settimanali, le decisioni tecniche, il Change Log e l'Issue Log. Gli SLA di risposta sono commisurati alla priorità, da poche ore per ciò che blocca il lavoro fino a qualche giorno per le richieste a bassa priorità, e una semplice etiquette (thread per le discussioni lunghe, tag solo quando serve un'azione) mantiene ordinato il flusso. I meeting sono riservati ai casi in cui servono davvero, secondo il calendario delle cerimonie definito nelle regole operative.

\section{Gestione del Cambiamento}

In presenza di una richiesta di modifica a scope, timeline o budget già approvati --- su iniziativa dello stakeholder, per un'impossibilità tecnica emersa in corso d'opera o per un evento esterno --- si attiva un processo di Change Request formalizzato in cinque passi. Alla submission del Change Request Form segue l'analisi di impatto: il Project Manager, insieme alla Tech Lead, redige un Project Impact Statement che quantifica l'effetto del cambiamento su scope, tempi, budget, qualità e risorse, e presenta le opzioni possibili (accettare estendendo la timeline, differire a una release successiva, rifiutare) con la raccomandazione motivata del PM. La decisione spetta allo sponsor e, in caso di impatto significativo su budget o timeline, comporta la rinegoziazione del contratto; seguono l'implementazione --- con l'aggiornamento di POS, WBS, Gantt e Product Backlog e l'assegnazione a un Responsible --- e il tracciamento nel Change Log fino alla chiusura. Questo presidio formale serve a evitare lo scope creep e a garantire decisioni consapevoli: i sottosistemi Agile accolgono i cambiamenti come parte naturale del processo di sviluppo, mentre per i sottosistemi tradizionali ogni modifica viene valutata caso per caso in termini di benefici, necessità e rischi. Per quanto riguarda l'allocazione delle risorse, in caso di conflitto con altri progetti di PlayHeritage Labs, MaraffaOnline ha la precedenza, evitando situazioni di over-allocation.
